PCCTP heuristics already combine construction, local improvement and
diversification. Cl\'{\i}maco et al.~\cite{climaco2021reactive} developed a
reactive GRASP using an adapted GENIUS construction and randomized variable
neighborhood descent. Their local search changes both route order and visited
locations while respecting the problem's feasibility requirements. This
provides a stronger point of comparison than a construction heuristic alone.
The experimental question here concerns the additional value of a separate
destruction-and-repair phase after initialization, rather than the first use
of multiple neighborhoods for the PCCTP.

Exact optimization and mathematical-programming hybrids also have direct
precedents for this problem. Cl\'{\i}maco et al.~\cite{climaco2023approaches}
presented a branch-and-cut algorithm, valid inequalities, reduction rules and
a hybrid combining a PCCTP heuristic with integer programming. Consequently,
neither adding preprocessing nor combining a GRASP-derived method with a MIP
solver is sufficient by itself to distinguish the present investigation.
Any useful extension must be specified more precisely and evaluated against
appropriate controls, including its model-building and optimization time.

The later data-mining approach of Cl\'{\i}maco et
al.~\cite{climaco2026mdm} exploits patterns in high-quality solutions to guide
GRASP search. Its publisher description includes analyses of pattern mining,
statistical comparisons and time-to-target behavior. The supplied MDM code
and manuscript support the recovered reference used here. We distinguish
that reference's declared corrections and parameter choices from the
published method description. In particular, a comparison ending before
mining is activated cannot establish superiority over the mining mechanism.
Final-cost comparisons and mining-exposure records address different parts
of that question.

Destruction and repair also belong to the broader large-neighborhood-search
literature. Ropke and Pisinger~\cite{ropke2006alns} used competing
subheuristics with selection frequencies reflecting their historical
performance in adaptive large neighborhood search for pickup and delivery
with time windows. The uniform-removal configuration evaluated here has no
such adaptive operator-selection layer. It should therefore not be described
as an implementation or improvement of that complete framework. Its
temporary coverage and prize violations, greedy reinsertion and acceptance
rule define the specific trajectory whose incremental value is examined.

The route-ordering component creates another attribution problem.
Helsgaun~\cite{helsgaun2000} developed an effective implementation of the
Lin--Kernighan traveling-salesman heuristic. In the present implementation,
LKH receives a symmetric routing matrix for a fixed set of visited locations.
Reordering that set cannot repair a missing coverage location or add prize;
those effects require visit changes. Conversely, an observed hybrid gain
could arise from strong reordering rather than the visit-change mechanism.
The fixed-support component replay is designed to separate these effects,
while retaining its conditional status and charged remaining allowance.

General-purpose MIP neighborhood methods offer related alternatives.
Fischetti and Lodi~\cite{fischetti2003local} used local branching cuts to
define solution subspaces explored by a MIP solver. Danna, Rothberg and
Le Pape~\cite{danna2005rins} introduced relaxation-induced neighborhood
search, which uses information from a continuous relaxation and formulates
neighborhood exploration as a further MIP. These methods motivate exact
exploration around feasible incumbents, but they do not identify a unique
PCCTP neighborhood. Fixing routing edges and visit decisions, contracting
retained paths, and restricting insertion positions impose distinct
structural choices. Their feasibility, computational cost and improvement
potential require separate evaluation; a local optimality certificate is
limited to the decisions left free.

This investigation therefore separates three empirical questions: whether
changing visits adds value beyond fixed-support reordering, whether restricted
integer repair adds value beyond greedy repair on matched destructions, and
whether any resulting complete search improves on a sufficiently developed
MDM reference under an explicit total-time comparison. The completed
exploratory tests include unfavorable outcomes and a failed validation
screen. The ongoing longer-horizon work does not yet settle the third
question. Differences in neighborhood construction are descriptions of the
method under study, not evidence of novelty or of a literature-wide
performance advantage.
